% 第 5 章 引用的重载（由 tools/convert.py 从原 content/ 自动生成，勿手改）

本章讨论不同引用限定符的重载成员函数。比如，关于getter应该通过值返回，还是通过常量引用返回的问题。\par


\section{getter的返回类型}

在C++11之前，当为复制开销很大的成员实现getter时，有以下几种选择:\par

\begin{itemize}
  \item 返回值
  \item 通过左值引用返回
\end{itemize}

简单地讨论一下这些方案。\par

\mySmallsection{5.1.1 传值返回}

按值返回的getter看起来是这样的(记住:不要使用\textit{const}按值返回，否则禁用移动语义):\par

\begin{codebox}[cpp]
class Person
{
private:
  std::string name;
public:
  ...
  std::string getName() const {
    return name;
  }
};
\end{codebox}

这段代码是安全的，但每次获取\textit{name}时，都可能复制\textit{name}。\par

例如，只是检查是否有没名字的人会有很大的开销:\par

\begin{codebox}[cpp]
std::vector<Person> coll;
...
for (const auto& person : coll) {
  if (person.getName().empty()) { // OOPS: copies the name
    std::cout << "found empty name\n";
  }
}
\end{codebox}

如果将此方法与返回引用的方法进行比较，可以看到按值返回字符串的版本的性能开销是引用的2到100倍(前提是名称的长度较大，这样SSO就没有帮助了)。对图像或数千个元素的集合的成员可能更糟糕。在这种情况下，getter通常返回(const)引用以提高性能。\par

\mySmallsection{5.1.2 引用返回}

通过引用返回的getter看起来像这样:\par

\begin{codebox}[cpp]
class Person
{
private:
  std::string name;
public:
  ...
  const std::string& getName() const {
    return name;
  }
};
\end{codebox}

这样更快，但不安全，因为调用者必须确保返回引用的对象生命周期足够长。实际上，使用返回引用值比原始对象的时间长的话，会有生命周期上的风险。\par

会踩入这种陷阱的一种方式是使用基于范围的for循环，如下所示:\par

\begin{codebox}[cpp]
for (char c : returnPersonByValue().getName()) { // OOPS: undefined behavior
  if (c == ' ') {
    ...
  }
}
\end{codebox}

注意，循环的右边有一个函数，它返回一个我们用getter引用的临时对象。但是，定义了基于范围的for循环，使上面的代码等价于以下代码:\par

\begin{codebox}[cpp]
reference range = returnPersonByValue().getName();
// OOPS: returned temporary object destroyed here
for (auto pos = range.begin(), end = range.end(); pos != end; ++pos) {
  char c = *pos;
  if (c == ' ') {
    ...
  }
}
\end{codebox}

开始迭代之前，初始化\textit{reference1}，必须使用传递的范围两次(一次调用begin()，一次调用end())，并希望避免创建副本(这可能代价很高，甚至不可能)。一般来说，引用会延长所引用内容的生命周期。在本例中，\textit{range}并不引用由\textit{returnPersonByValue(\allowbreak{})\allowbreak{}}返回的\textit{Person}，而\textit{range}指向\textit{getName()}的返回值，是对返回的\textit{Person}的引用。因此，\textit{range}扩展了引用的生命周期。因此，在第一个语句的末尾，返回的临时对象将销毁，并且在遍历名称的字符时使用对已销毁对象名称的引用。\par

最好的情况下，在这里会得到一个核心转储(段错误)，这样就可以看到出现了明显的错误。最坏的情况是，发布了软件，并出现了未定义行为。\par

如果getter按值返回名称，这样的代码就不会有问题。这样，\textit{range}将扩展名称副本的生存期，以便我们使用该名称直到\textit{range}的生存期结束。\par

\mySmallsection{5.1.3 使用移动语义解决困境}

通过移动语义，就可以解决这个困境的方法。如果这样做安全，可以通过引用返回。如果遇到生命周期的问题，可以通过值返回。\par

方法如下:\par

\begin{codebox}[cpp]
class Person
{
private:
  std::string name;
public:
  ...
  std::string getName() && { // when we no longer need the value
    return std::move(name); // we steal and return by value
  }
  const std::string& getName() const& { // in all other cases
    return name; // we give access to the member
  }
};
\end{codebox}

用不同的引用限定符重载getter，方法与重载带有\&\&和\textit{const}\&形参的函数相同:\par

\begin{itemize}
  \item 带有\&\&限定符的是有不再需要值的对象时使用(一个即将死亡的对象或用\textit{std::move()}标记的对象)。
  \item 带有\textit{const}\&限定符的版本用于所有其他情况。总是合适的，但只是无法获得\&\&版本时的备选方案。因此，如果有一个不会销毁的对象或没标记为\textit{std::move()}，就会使用此函数。
\end{itemize}

现在的性能和安全性都很好:\par

\begin{codebox}[cpp]
Person p{"Ben"};
std::cout << p.getName(); // 1) fast (returns reference)
std::cout << returnPersonByValue().getName(); // 2) fast (uses move())

std::vector<Person> coll;
...
for (const auto& person : coll) {
  if (person.getName().empty()) { // 3) fast (returns reference)
    std::cout << "found empty name\n";
  }
}

for (char c : returnPersonByValue().getName()) { // 4) safe and fast (uses move())
  if (c == ' ') {
    ...
  }
}
\end{codebox}

语句1)和3)使用\textit{const}\&的版本，因为有一个对象没有使用\textit{std::move()}标记。语句2)和4)使用\&\&的版本，因为我们为一个临时对象调用\textit{getName()}。因为临时对象即将销毁，getter可以将成员名移出作为返回值，这意味着不必为返回值分配新的内存，可以窃取了成员的值。\par

您可能还记得，return语句不应该使用\textit{std::move()}来返回已经销毁的局部对象。本例中，不返回任何局部对象，而是返回一个成员，其生存期不以成员函数的结束为结束。\par

\mySmallsection{\textit{std::move()}用于成员函数中}

请注意，这个特性意味着即使在调用成员函数时也可以使用\textit{std::move()}。例如:\par

\begin{codebox}[cpp]
void foo()
{
  Person p{ ... };
  ...
  coll.push_back(p.getName()); // calls getName() const&
  ...
  coll.push_back(std::move(p).getName()); // calls getName() && (OK, p no longer used)
}
\end{codebox}

使用\textit{getName()}时，使用\textit{std::move()}将提高程序的性能。不返回对\textit{const} std::string的引用(只能复制)，返回值是作为非\textit{const}字符串返回的\textit{p}的移动名，因此\textit{push\_\allowbreak{}back(\allowbreak{})\allowbreak{}}可以使用移动语义将其移动到\textit{coll}中。通常，在这之后，\textit{p}处于有效但未定义的状态。\par

关于在C++标准库中使用此特性的示例，请参见std::optional<>。\par


\section{重载}

从C++98开始，可以重载成员函数来实现const和非const版本。例如:\par

\begin{codebox}[cpp]
class C {
public:
  ...
  void foo(); // foo() for non-const objects
  void foo() const; // foo() for const objects
};
\end{codebox}

圆括号后面的限定符允许限定一个没有传递给形参的对象:可以调用此成员函数的对象。\par

有了移动语义，就有了用限定符重载函数的新方。考虑以下代码:\par

\begin{codeboxt}[cpp]{basics/refqual.cpp}
#include <iostream>
class C {
public:
  void foo() const& {
    std::cout << "foo() const&\n";
  }
  void foo() && {
    std::cout << "foo() &&\n";
  }
  void foo() & {
    std::cout << "foo() &\n";
  }
  void foo() const&& {
    std::cout << "foo() const&&\n";
  }
};

int main()
{
  C x;
  x.foo(); // calls foo() &
  C{}.foo(); // calls foo() &&
  std::move(x).foo(); // calls foo() &&

  const C cx;
  cx.foo(); // calls foo() const&
  std::move(cx).foo(); // calls foo() const&&
}
\end{codeboxt}

这个程序演示了所有可能的引用限定符，以及何时调用。通常，只有两到三个这样的重载，比如对getter使用\&\&和const\&(和\&)。\par

还要注意，不允许引用和非引用限定符的重载:\par

\begin{codebox}[cpp]
class C {
public:
  void foo() &&;
  void foo() const; // ERROR: can’t overload by both reference and value qualifiers
};
\end{codebox}


\section{何时使用引用}

引用限定符允许对特定类别的对象调用函数时，以不同的方式实现函数。目标是在为不再需要其值的对象，调用不同的成员函数。\par

虽然确实有这个特性，但没有使用。可以(也应该)用它来确保修改对象的操作，不会让即将销毁的临时对象调用。\par

\mySmallsection{5.3.1 赋值操作符的引用限定符}

更好地使用引用限定符的方式是修改赋值操作符的实现。如\texthref{http://wg21.link/n2819}中建议的那样，在可能的地方使用引用限定符声明赋值操作符可能会更好。\par

例如，字符串的赋值操作符声明如下:\par

\begin{codebox}[cpp]
namespace std {
  template<typename charT, ... >
  class basic_string {
  public:
    ...
    constexpr basic_string& operator=(const basic_string& str);
    constexpr basic_string& operator=(basic_string&& str) noexcept( ... );
    constexpr basic_string& operator=(const charT* s);
    ...
  };
}
\end{codebox}

这允许将新值赋给临时字符串:\par

\begin{codebox}[cpp]
std::string getString();

getString() = "hello"; // OK
foo(getString() = ""); // passes string instead of bool
\end{codebox}

考虑用引用限定符声明赋值操作符:\par

\begin{codebox}[cpp]
namespace std {
  template<typename charT, ... >
  class basic_string {
  public:
    ...
    constexpr basic_string& operator=(const basic_string& str) &;
    constexpr basic_string& operator=(basic_string&& str) & noexcept( ... );
    constexpr basic_string& operator=(const charT* s) &;
    ...
  };
}
\end{codebox}

代码将不再编译:\par

\begin{codebox}[cpp]
std::string getString();

getString() = "hello"; // ERROR
foo(getString() = ""); // ERROR
\end{codebox}

注意，特别是对于可以用作布尔值的类型，将有助于找到如下错误:\par

\begin{codebox}[cpp]
std::optional<int> getValue();

if (getValue() = 0) { // OOPS: compiles although = is used instead of ==
  ...
}
\end{codebox}

其实，会给临时对象返回基本类型的属性:右值。\par

注意，所有这些修改C++标准的建议都被拒绝了，主要是考虑向后兼容性。然而，在实现自己的类时，可以使用以下改进的方式:\par

\begin{codebox}[cpp]
class MyType {
public:
  ...
  // disable assigning value to temporary objects:
  MyType& operator=(const MyType& str) & =default;
  MyType& operator=(MyType&& str) & =default;

  // because this disables the copy/move constructor, also:
  MyType(const MyType&) =default;
  MyType(MyType&&) =default;
  ...
};
\end{codebox}

通常，对可能修改对象的成员函数执行此操作。\par

\mySmallsection{5.3.2 其他成员函数的引用限定符}

如getter示例所示，引用限定符也可以而且应该在返回对对象的引用时使用。这样，可以减少访问已销毁临时对象的成员的风险。\par

同样，标准字符串的当前声明可以作为一个例子:\par

\begin{codebox}[cpp]
namespace std {
  template<typename charT, ... >
  class basic_string {
  public:
    ...
    constexpr const charT& operator[](size_type pos) const;
    constexpr charT& operator[](size_type pos);
    constexpr const charT& at(size_type n) const;
    constexpr charT& at(size_type n);

    constexpr const charT& front() const;
    constexpr charT& front();
    constexpr const charT& back() const;
    constexpr charT& back();
    ...
  };
}
\end{codebox}

相反，下面的重载会更好一些:\par

\begin{codebox}[cpp]
namespace std {
  template<typename charT, ... >
  class basic_string {
  public:
    ...
    constexpr const charT& operator[](size_type pos) const&;
    constexpr charT& operator[](size_type pos) &;
    constexpr charT operator[](size_type pos) &&;
    constexpr const charT& at(size_type n) const&;
    constexpr charT& at(size_type n) &;
    constexpr charT at(size_type n) &&;

    constexpr const charT& front() const&;
    constexpr charT& front() &;
    constexpr charT front() &&;
    constexpr const charT& back() const&;
    constexpr charT& back() &;
    constexpr charT back() &&;
    ...
  };
}
\end{codebox}

同样，由于向后兼容性，C++标准中相应的更改可能会成为一个问题。但可以为自定义的类型提供这些重载。这种情况下，不要忘记右值引用的实现可以移出大型成员。\par


\section{总结}

\begin{itemize}
  \item 可以在不同的引用限定符上重载成员函数。
  \item 重载getter用于具有引用限定符的大型成员，使其既安全又快速。
  \item 使用\textit{std::move()}标记对象是有意义的，即使在调用成员函数时也是如此。
  \item 在赋值操作符中使用引用限定符。
\end{itemize}
